\section{Conclusiones, Implicancias y Reproducibilidad}

% ================================================================
\begin{frame}{Conclusiones e Implicancias de Política Regional}
\vspace{-0.18cm}
\begin{columns}[T,onlytextwidth]
\begin{column}{0.48\textwidth}
\begin{contrastbox}{ccdenavy}{Hallazgos Empíricos Centrales}
\scriptsize\RaggedRight
\begin{enumerate}\setlength{\itemsep}{0.06em}
  \item \textbf{Jerarquía dual ($K=4 \subset K=2$):} CLR+PCA2 divide al país en un bloque Minero ($N=10$) y uno Diversificado ($N=14$), subdivididos con $85\%$ de estabilidad bootstrap en 4 perfiles.
  \item \textbf{Minería Diversificada vs.\ Enclave:} \textit{Minero/Div.} (Arequipa, Moquegua, Cusco) superó su nivel pre-COVID ($106.6$ en 2023), pero \textit{Minero/Enclave} (Apurímac, Pasco, Huancavelica) sufrió una cicatriz de $-10.92$ pts ($p=0.0029$).
  \item \textbf{Trampa de volatilidad:} Los enclaves soportan una volatilidad ($\sigma=20.5\%$) $4.4\times$ mayor que el grupo agro/servicios ($\sigma=4.6\%$).
\end{enumerate}
\end{contrastbox}
\end{column}

\begin{column}{0.48\textwidth}
\begin{contrastbox}{okgreen}{Agenda de Política Regional}
\scriptsize\RaggedRight
\begin{itemize}\setlength{\itemsep}{0.04em}
  \item \textbf{Encadenamientos en enclaves:} Invertir canon en Apurímac, Pasco y Huancavelica en agroindustria y logística que amortigüen paros extractivos.
  \item \textbf{Fondos de estabilización:} Reglas contracíclicas en clusters con $\sigma > 15\%$ para suavizar inversión pública.
  \item \textbf{Métrica en niveles:} Evaluar desempeño con índices acumulados ($\text{PIB}_{2019=100}$) en vez de tasas anuales.
\end{itemize}
\end{contrastbox}
\vspace{0.03cm}
\begin{keybox}[left=2mm,right=2mm,top=0.6mm,bottom=0.6mm]
\scriptsize\RaggedRight \textbf{Mensaje Final}: La fragilidad pospandemia proviene de la \textbf{falta de diversificación} en economías de enclave.
\end{keybox}
\end{column}
\end{columns}
\source{Síntesis del equipo de investigación -- Círculo de Ciencia de Datos y Econometría (CCDE).}
\end{frame}

% ================================================================
\begin{frame}{Trazabilidad Reproducible y Referencias Clave}
\vspace{-0.14cm}
\begin{columns}[T,onlytextwidth]
\begin{column}{0.51\textwidth}
\begin{ccdebox}[left=2mm,right=2mm,top=1mm,bottom=1mm]
\scriptsize
\textbf{\textcolor{ccdenavy}{Arquitectura del Repositorio}}
\vspace{0.04cm}

\resizebox{\linewidth}{!}{%
\renewcommand{\arraystretch}{1.10}
\begin{tabular}{@{}l l@{}}
\toprule
\textbf{Módulo / Archivo} & \textbf{Función en el Pipeline} \\
\midrule
\texttt{src/generador-df.py} & ETL de 24 Excels INEI $\rightarrow$ panel base \\
\texttt{src/clustering\_paso1..4.py} & Ceros $\delta=0.05$, CLR/ILR, PCA y Gap \\
\texttt{src/econ\_paso1..3.py} & Modelos en tasas y diagnóstico $K=2$ \\
\texttt{src/generar\_analisis\_completo.py} & Modelos M1--M6 en niveles y 6 figuras \\
\texttt{src/modelo-clusters.ipynb} & Cuaderno ejecutable de tipología CoDa \\
\texttt{src/modelo-econometrico.ipynb} & Cuaderno ejecutable de Panel FE \\
\texttt{data/processed/*.xlsx} & Bases (\texttt{panel\_econometrico\_final}) \\
\bottomrule
\end{tabular}%
}
\end{ccdebox}
\vspace{0.05cm}
\begin{keybox}[left=2mm,right=2mm,top=0.9mm,bottom=0.9mm]
\scriptsize\RaggedRight \textbf{Reproducibilidad}: Semilla fija (\texttt{RANDOM\_STATE = 0}), entorno declarativo en \texttt{requirements.txt} y trazabilidad en Git.
\end{keybox}
\end{column}

\begin{column}{0.46\textwidth}
\begin{ccdebox}[left=2mm,right=2mm,top=1mm,bottom=1mm]
\scriptsize\RaggedRight
\textbf{\textcolor{ccdenavy}{Referencias Metodológicas Clave}}
\vspace{0.05cm}

\textbf{Aitchison, J.} \textcolor{ccdemid}{(1986)}. \textit{The Statistical Analysis of Compositional Data}. Chapman \& Hall.\\[0.06cm]
\textbf{Martín-Fernández, J. A. et al.} \textcolor{ccdemid}{(2003)}. Dealing with zeros in compositional data sets. \textit{Math.\ Geol.}, 35(3), 253--278.\\[0.06cm]
\textbf{Tibshirani, R., Walther, G. \& Hastie, T.} \textcolor{ccdemid}{(2001)}. Estimating the number of clusters via the gap statistic. \textit{JRSS-B}, 63(2), 411--423.\\[0.06cm]
\textbf{Wooldridge, J. M.} \textcolor{ccdemid}{(2010)}. \textit{Econometric Analysis of Cross Section and Panel Data}. MIT Press.\\[0.06cm]
\textbf{INEI} \textcolor{ccdemid}{(2025)}. \textit{Producto Bruto Interno por Departamentos 2007--2025}. Lima.
\end{ccdebox}
\end{column}
\end{columns}
\source{Repositorio oficial: \texttt{https://github.com/CCDE-DataLab/clusters-departamentales}.}
\end{frame}